\documentclass{article}
\usepackage{amsmath,amssymb}
\usepackage{xurl}
\usepackage[hidelinks]{hyperref}
\setlength{\emergencystretch}{2em}
% Shared commands for notes in docs/paper-gaps/.

\DeclareUnicodeCharacter{2080}{\ifmmode{}_{0}\else\textsubscript{0}\fi}
\DeclareUnicodeCharacter{2081}{\ifmmode{}_{1}\else\textsubscript{1}\fi}
\DeclareUnicodeCharacter{2082}{\ifmmode{}_{2}\else\textsubscript{2}\fi}
\DeclareUnicodeCharacter{2099}{\ifmmode{}_{n}\else\textsubscript{n}\fi}

\newcommand{\Lean}{\textsc{Lean}}
\ExplSyntaxOn
\NewDocumentCommand{\leanid}{m}{
  \ifmmode
    \textsf{\detokenize{#1}}
  \else
    \group_begin:
    \urlstyle{sf}
    \tl_set:Nn \l_tmpa_tl {#1}
    \tl_replace_all:Nnn \l_tmpa_tl {\_} {_}
    \exp_args:NV \path \l_tmpa_tl
    \group_end:
  \fi
}
\ExplSyntaxOff
\providecommand{\tr}{\operatorname{tr}}
\newcommand{\tnleanIssueBase}{https://github.com/LionSR/TNLean/issues/}
\newcommand{\tnleanPrBase}{https://github.com/LionSR/TNLean/pull/}
\newcommand{\tnleanDocBase}{https://sirui-lu.com/TNLean/docs/}
\newcommand{\ghissue}[1]{\href{\tnleanIssueBase#1}{GitHub issue \##1}}
\newcommand{\ghpr}[1]{\href{\tnleanPrBase#1}{PR \##1}}
\newcommand{\doclink}[1]{\href{\tnleanDocBase#1.html}{\path{#1}}}

% Dirac notation and the identity operator, as in blueprint/src/macros/common.tex.
% \providecommand keeps note-local definitions authoritative.
\usepackage{dsfont}
\providecommand{\ket}[1]{|#1\rangle}
\providecommand{\bra}[1]{\langle #1|}
\providecommand{\braket}[2]{\langle #1 | #2 \rangle}
\providecommand{\ketbra}[2]{|#1\rangle\!\langle #2|}
\providecommand{\Id}{\mathds{1}}
\providecommand{\one}{\mathds{1}}
\providecommand{\C}{\mathbb{C}}
\providecommand{\MN}[1]{M_{#1}(\C)}
\providecommand{\MD}{M_D(\C)}

% Verdict marker: \gapnote{<kind>}{<status>}. Kind is the policy
% classification (clarification, local-correction, scope-restriction,
% unfaithful, false-source, open-gap); status is open, wip (elimination
% underway), resolved, or historical. Machine-read by the site index;
% prints a verdict line.
\newcommand{\gapnote}[2]{\par\noindent\textbf{Verdict:} #1 (#2).\par}

\title{The Mixed Interpolation Path: One-Site Injectivity of the Common Operator}
\author{}
\date{2026-11-17}
\begin{document}
\maketitle
\gapnote{false-source}{resolved}

\section{The source claim}
Garre-Rubio, Lootens and Moln\'ar connect two injective matrix product
states $A_0,A_1$, symmetric under two matrix product operator
representations of one fusion category $\mathcal C$ with equal $F$-symbols
and equal $L$-symbols, by the interpolating tensor $A(\gamma)$ and a common
operator $T_a$ for every simple object $a$
\cite[Section~5]{GarreRubio2023MPOSym}. The operator $T_a$ keeps the endpoint
operators $T^0_a,T^1_a$ in the two diagonal physical sectors and contracts
the action tensors of the two endpoints in the two crossed sectors. The
source then states:
\begin{quote}
    ``Let us first show that $T_a$ is injective. Since we $T_a^0$ and
    $T_a^1$ are injective we only need to show that
    $T_a^{0,1} \oplus T_a^{1,0}$ are injective or equivalently, that there
    is a left inverse acting on the physical level. For that we map
    $T_a^{0,1} \oplus T_a^{1,0}$ to
    $A_0\otimes A_1( T_a^{0,1} \oplus T_a^{1,0})
        A_0^{-1}\otimes A_1^{-1}$
    under an appropriate rearrangement of the indices. This results in
    $T_a^{0} A_0A_0^{-1} \oplus T_a^{1} A_1A_1^{-1}$ which is
    invertible since we have restricted to the subspace generated by the
    tensors so that $A_pA_p^{-1}= \mathbb 1_{\mathcal{H}_p}$.''
\end{quote}
The block-injective case repeats the claim: ``The injectivity of every block
of $A(\gamma)$ and the injectivity of $T_a$ is justified in the same way as
in the previous case.''\footnote{Source traceability:
    \path{Papers/2203.12563/REsubmission.tex}, lines 1584--1631 (the tensor
    $A(\gamma)$ and the operator $T_a$), lines 1632--1685 (the action and
    the alignment of $F$ and $L$), lines 1687--1689 (the claim quoted
    above), and lines 1695 onwards (the block-injective case).}
Here injectivity is one-site injectivity: the left inverse acts on one
physical site.

This note shows that the claim is false. The second claim of the same
paragraph, injectivity of $A(\gamma)$ for $0<\gamma<1$, is true with the
source's hypotheses and is formalized.\footnote{%
    \leanid{MPSTensor.MPOSymmetry.isInjective_mixedEndpointInterpolation},
    in \doclink{TNLean/MPS/Symmetry/MPOSymmetry/MixedEndpointInterpolation}.}

\section{The letters of the common operator}
Restricting each local space to the span of $A_p$, as the source does
(lines 1584--1586), the $D_p^2$ matrices $A_p^{ij}$ form a basis of
$\mathcal M_{D_p}$, and after a change of physical basis they are the matrix
units $|i\rangle\langle j|$. Write $E_p=\mathbb C^{D_p}$, $K_p=\mathbb C^{\chi_p}$
for the bond space of $T^p_a$, and $m$ for the common action
multiplicity of the object $a$. The exact action of $T^p_a$ on $A_p$ with
biorthogonal action maps $W_{p,\mu}:E_p\to K_p\otimes E_p$ and
$V_{p,\mu}:K_p\otimes E_p\to E_p$, $V_{p,\mu}W_{p,\nu}=\delta_{\mu\nu}\mathbb 1$,
then determines the endpoint letters
\begin{align}
    (T^p_a)^{(r,s),(k,l)}_{\alpha\beta}
     & =\sum_{\mu=0}^{m-1}(W_{p,\mu})_{(\alpha,k),r}(V_{p,\mu})_{s,(\beta,l)}.
    \label{eq:glm23_oneSite_endpoint}
\end{align}
The crossed letters of $T_a$ are given by the same contraction with the
analysis map of the other endpoint, $p\ne q$:
\begin{align}
    (T_a)^{(r,s),(k,l)}_{\alpha\beta}
     & =\sum_{\mu=0}^{m-1}(W_{p,\mu})_{(\alpha,k),r}(V_{q,\mu})_{s,(\beta,l)},
     &
    r,k\in E_p,\ s,l\in E_q,\ \alpha\in K_p,\ \beta\in K_q.
    \label{eq:glm23_oneSite_crossed}
\end{align}
Each letter of $T_a$ is supported on one of the four corners
$K_p\times K_q$ of $K_0\oplus K_1$. Hence $T_a$ is one-site injective if and
only if each corner is spanned by its letters. The diagonal corners are
spanned because $T^0_a,T^1_a$ are injective, so the claim is equivalent to
the two crossed corners being spanned.\footnote{The formal operator is
    \leanid{MPSTensor.MPOSymmetry.mixedEndpointMPO}, in
    \doclink{TNLean/MPS/Symmetry/MPOSymmetry/MixedEndpointMPO}; its crossed
    letters are (\ref{eq:glm23_oneSite_crossed}).}

\paragraph{Multiplicity one.} If $m=1$ the claim holds. Write
$w_{kr}\in K_p$ and $v_{sl}\in K_q$ for the vectors
$\alpha\mapsto(W_p)_{(\alpha,k),r}$ and $\beta\mapsto(V_q)_{s,(\beta,l)}$.
The crossed letters are $w_{kr}v_{sl}^{\mathsf T}$, so their span is
$\operatorname{span}\{w_{kr}\}\otimes\operatorname{span}\{v_{sl}\}$. The
diagonal letters of $T^p_a$ are $w_{kr}$ times the analogous vectors of
$V_p$, so injectivity of $T^p_a$ forces $\operatorname{span}\{w_{kr}\}=K_p$,
and likewise $\operatorname{span}\{v_{sl}\}=K_q$. In particular the claim
holds for every group, where every action multiplicity is one. For $m\ge2$
the contraction over $\mu$ in (\ref{eq:glm23_oneSite_crossed}) can lose
rank, and it does in the following example.

\section{A counterexample}
Let $\mathcal C=\operatorname{Rep}(A_4)$, whose simple objects are the three
characters of $A_4$ and the three-dimensional irreducible representation
$\rho$, obtained from the permutation action on $\{0,1,2,3\}$ restricted to
vectors with coordinate sum zero. Let $c=(123)$, $x=(01)(23)$, $y=(02)(13)$
and $c'=(032)$, and take the label sets
\begin{align}
    S_0 & =\{e,c,x,y\}, & S_1 & =\{e,c,x,c'\}.
    \label{eq:glm23_oneSite_labels}
\end{align}
For $p\in\{0,1\}$ let $E_p=\mathbb C^{S_p}$, so $D_p=4$, let $A_p$ be the
sixteen matrix units of $\mathcal M_4$, and, for every irreducible unitary
representation $\rho_a$ of $A_4$ of dimension $d_a$, set
$K_p=\mathbb C^{d_a}$, $m=d_a$ and
\begin{align}
    W_{p,\mu}|r\rangle & =\rho_a(r)^{-1}e_\mu\otimes|r\rangle,
                       &
    V_{p,\mu}(\kappa\otimes|s\rangle)
                       & =(\rho_a(s)\kappa)_\mu\,|s\rangle,
    \qquad r,s\in S_p.
    \label{eq:glm23_oneSite_action}
\end{align}
Then $V_{p,\mu}W_{p,\nu}=\delta_{\mu\nu}\mathbb 1$, and
(\ref{eq:glm23_oneSite_endpoint}) gives the diagonal operator whose letter on
the physical pair $(r,s)$ is $\rho_a(r^{-1}s)$ and which vanishes off the
diagonal of physical indices. Every hypothesis of the source holds.
\begin{enumerate}
    \item \emph{States.} Each $A_p$ is injective, with physical space
          $E_p\otimes E_p$ equal to the span of the tensor.
    \item \emph{Endpoint operators.} A one-dimensional $\rho_a$ gives a
          bond-one operator with nonzero letters, which is injective; the
          trivial object gives the identity. For $\rho$, the set
          $S_p^{-1}S_p$ has $10$ elements for $p=0$ and $11$ for $p=1$,
          and $\rho(S_p^{-1}S_p)$ spans $\mathcal M_3$. So every $T^p_a$ is
          injective.
    \item \emph{Fusion and $F$-symbols.} Fix Clebsch--Gordan isometries
          $U^{c,\nu}_{ab}:\mathbb C^{d_c}\to\mathbb C^{d_a}\otimes\mathbb C^{d_b}$
          of $A_4$. The stacked letter of $T^p_aT^p_b$ on $(r,s)$ is
          $(\rho_a\otimes\rho_b)(r^{-1}s)
              =\sum_{c,\nu}U^{c,\nu}_{ab}\rho_c(r^{-1}s)U^{c,\nu\,*}_{ab}$,
          so the maps $U^{c,\nu}_{ab}$ and their adjoints are exact fusion
          tensors of both endpoint algebras, for arbitrary boundaries. Both
          endpoints use the same fusion tensors, which determine the
          $F$-symbols, so $F^0=F^1$: the Wigner $6j$-symbols of $A_4$.
    \item \emph{Action and $L$-symbols.} The action
          (\ref{eq:glm23_oneSite_action}) is exact and biorthogonal with the
          same multiplicity $m=d_a$ at both endpoints. Applying
          $U^{c,\nu\,*}_{ab}$ to the composite action of $a$ and $b$ gives
          $U^{c,\nu\,*}_{ab}(\rho_a(r)^{-1}\otimes\rho_b(r)^{-1})
              =\rho_c(r)^{-1}U^{c,\nu\,*}_{ab}$, and similarly for the analysis
          maps. The $L$-symbols are therefore the Clebsch--Gordan
          coefficients at both endpoints, independently of $r$, so
          $L^0=L^1$ exactly.
\end{enumerate}
For $a=\rho$ the crossed letters (\ref{eq:glm23_oneSite_crossed}) of the
corner $K_0\times K_1$ are $\rho(g^{-1}h)$ with $g\in S_0$, $h\in S_1$,
and those of the other corner are their inverses. Let
\begin{align}
    M          & =\begin{pmatrix}
                      0  & 0  & 1  & -1 \\
                      1  & -1 & 0  & 0  \\
                      -1 & 1  & 0  & 0  \\
                      0  & 0  & -1 & 1
                  \end{pmatrix},
               &
    \varphi(q) & =\sum_{j=0}^{3}M_{q(j),j}=\operatorname{tr}(M^{\mathsf T}P_q),
    \label{eq:glm23_oneSite_coefficient}
\end{align}
where $P_qe_j=e_{q(j)}$ is the permutation matrix of $q$, and products
are composites, $(pq)(j)=p(q(j))$. The rows and columns of $M$
sum to zero, so $\varphi$ is a matrix coefficient of $\rho$:
$\varphi(q)=\ell(\rho(q))$ for a linear functional $\ell$ on $\mathcal M_3$.
The set $S_0^{-1}S_1$ consists of ten elements, and $\varphi$ vanishes on
each of them, while $\varphi$ takes the values $4$ and $-4$ at the two
remaining elements $(012)$ and $(023)$ of $A_4$. Thus $\ell\ne0$
annihilates every crossed letter of the corner $K_0\times K_1$, and the
corner is not spanned. The letters of the other corner are the inverses
$\rho(h^{-1}g)$, and $q\mapsto\varphi(q^{-1})$ is again a nonzero matrix
coefficient of $\rho$ vanishing on them, so that corner is not spanned
either. So $T_\rho$
is not one-site injective: its letters span a subspace of dimension $34$ of
the $36$-dimensional algebra $\mathcal M_6$.

\paragraph{Gauge choices.} The source fixes the action tensors only up to
the gauge freedom that preserves the $L$-symbols. Given $T^p_a$ and the
injective $A_p$, the biorthogonal action maps are unique up to an
invertible change $g_a$ of the multiplicity space: for two such families,
$V'_{p,\nu}W_{p,\mu}$ commutes with every $A_p^{ij}$ and is therefore a
scalar. The changes that preserve the $L$-symbols, that is, the families
with $U^{c,\nu\,*}_{ab}(g_a\otimes g_b)=g_cU^{c,\nu\,*}_{ab}$, are the tensor
automorphisms of the forgetful functor of $\operatorname{Rep}(A_4)$, which
by Tannaka duality are $g_a=\rho_a(z)$ with $z\in A_4$. Replacing
$W_{1,\mu},V_{1,\mu}$ by their transforms under $\rho_a(z)$ leaves
$T^1_a$ unchanged and replaces the crossed letters by $\rho(g^{-1}zh)$. For
every one of the twelve elements $z$ both crossed corners still span a
subspace of dimension $8$. No admissible choice of action tensors makes
$T_\rho$ injective.

\paragraph{Exact verification.} The script
\path{scripts/check_glm23_mixed_operator_one_site_injectivity.py} checks,
in exact rational arithmetic and in the index convention of the formal
operator, that $\rho$ is a representation, the biorthogonality and the
exact action at both endpoints, the injectivity of both endpoint operators
for $\rho$, the ranks $8$ of both crossed corners and $34$ of the full letter
span, the same ranks for all twelve gauges $z$, and the vanishing of
$\varphi$ on $S_0^{-1}S_1$. It uses an integral basis of $\rho$; a change
of bond basis conjugates every letter and preserves all ranks.

\section{Why this is a false-source claim}
The data are nondegenerate: both states are injective with nonzero
letters, every endpoint operator is injective, the multiplicities are the
positive dimensions $d_a$, and the trivial object acts as the identity.
No reading of the source's definitions excludes the example, so the
failure is not a degenerate reading. The source's argument replaces the
crossed corner by $T^0_aA_0A_0^{-1}\oplus T^1_aA_1A_1^{-1}$; but the
crossed letters (\ref{eq:glm23_oneSite_crossed}) contract the synthesis of
one endpoint with the analysis of the other through the multiplicity
index, and that contraction is not a rearrangement of either endpoint
operator. With multiplicity one the contraction is an outer product and the
claim survives; with multiplicity at least two it can lose rank. The
block-injective case inherits the counterexample, since one block is an
instance of it.

\section{The formal substitute}
The common operator is injective after blocking three sites, hence
normal, from fewer hypotheses: injective endpoint operators and
biorthogonal action maps with positive common multiplicity; positivity
follows from an exact action on an injective state.\footnote{%
    \leanid{MPSTensor.MPOSymmetry.isNBlkInjective_mixedEndpointMPO_three},
    \leanid{MPSTensor.MPOSymmetry.isNormal_mixedEndpointMPO} and
    \leanid{MPSTensor.MPOSymmetry.isNBlkInjective_mixedEndpointMPO_three_of_exact_action},
    in \doclink{TNLean/MPS/Symmetry/MPOSymmetry/MixedEndpointMPONormality}.}
The one-site claim is not formalized,
and should not be: it is false. The positive multiplicity-one statement
above is not formalized either.

\bibliographystyle{alpha}
\bibliography{references}
\end{document}
